% Part: incompleteness
% Chapter: arithmetization-syntax
% Section: introduction

\documentclass[../../../include/open-logic-section]{subfiles}

\begin{document}

\olfileid[id]{inc}{art}{int}
\olsection{Pendahuluan}

Untuk menghubungkan komputabilitas dengan logika, kita memerlukan cara untuk
membicarakan objek-objek logika (simbol, suku, !!{formula}s,
!!{derivation}s), operasi atas objek-objek itu, serta sifat dan relasi di
antara mereka, dengan cara yang dapat ditangani secara komputasional. Kita
dapat melakukannya secara langsung dengan meninjau fungsi dan relasi yang dapat
dihitung pada simbol, barisan simbol, dan objek lain yang dibangun darinya.
Karena semua objek sintaksis logis berhingga dan dibangun dari himpunan simbol
yang !!a{enumerable}, hal ini mungkin dilakukan untuk beberapa model
komputasi. Namun, model komputasi lain---seperti fungsi rekursif---terbatas
pada bilangan serta relasi dan fungsi atas bilangan. Selain itu, pada akhirnya
kita juga ingin dapat menangani sintaksis di dalam teori-teori tertentu,
khususnya teori yang dirumuskan dalam bahasa aritmetika. Dalam kasus-kasus ini,
kita perlu \emph{mengaritmetisasi} sintaksis, yakni merepresentasikan objek
sintaksis, operasi atasnya, dan relasi di antaranya masing-masing sebagai
bilangan, fungsi aritmetis, dan relasi aritmetis. Gagasan yang dapat dilacak
kembali sampai Leibniz ini ialah memasangkan bilangan dengan objek sintaksis.

Memasangkan bilangan sebagai ``kode'' bagi simbol relatif mudah. Beberapa
simbol sedikit menantang karena, misalnya, terdapat tak berhingga banyak
!!{variable}s, bahkan tak berhingga banyak !!{function}s untuk setiap
aritas~$n$. Namun, tentu saja kita dapat memasangkan bilangan dengan simbol
secara sistematis sedemikian rupa sehingga, misalnya, $\Obj v_2$ dan
$\Obj v_3$ memperoleh kode yang berbeda. Barisan simbol (seperti suku dan
!!{formula}s) lebih menantang. Akan tetapi, jika kita dapat menangani barisan
bilangan secara murni aritmetis (misalnya, melalui pengodean barisan dengan
pangkat bilangan prima), kita dapat memperluas pengodean simbol tunggal
menjadi pengodean barisan simbol, lalu memperluasnya lagi ke barisan atau
susunan !!{formula}s lainnya, seperti !!{derivation}s. Pengodean yang
diperluas ini disebut ``penomoran G\"odel.'' Setiap suku, !!{formula}, dan
!!{derivation} dipasangkan dengan sebuah bilangan G\"odel.

Dengan mengodekan barisan simbol sebagai barisan kode simbol-simbolnya dan
memilih sistem pengodean barisan yang dapat ditangani menggunakan fungsi dapat
dihitung, kita kemudian juga dapat menangani bilangan G\"odel menggunakan
fungsi dapat dihitung. Dalam praktiknya, semua fungsi yang relevan akan
rekursif primitif. Sebagai contoh, menghitung panjang suatu barisan dan
menghitung elemen ke-$i$ suatu barisan dari kode barisan tersebut keduanya
rekursif primitif. Jika bilangan yang mengodekan barisan itu, misalnya,
merupakan bilangan G\"odel !!a{formula}~$!A$, kita segera melihat bahwa
panjang !!a{formula} dan (kode) simbol ke-$i$ dalam !!a{formula} juga dapat
dihitung dari bilangan G\"odel~$!A$. Sedikit lebih sulit untuk membuktikan
bahwa, misalnya, sifat menjadi bilangan G\"odel suatu suku yang terbentuk
dengan benar atau suatu !!{derivation} yang benar bersifat rekursif primitif.
Meskipun demikian, hal itu mungkin karena barisan yang kita tinjau (suku,
!!{formula}s, !!{derivation}s) didefinisikan secara induktif.

Sebagai contoh, tinjaulah operasi substitusi. Jika $!A$ adalah formula, $x$
variabel, dan $t$ suku, maka $\Subst{!A}{t}{x}$ adalah hasil penggantian setiap
kemunculan bebas~$x$ dalam~$!A$ dengan~$t$. Sekarang andaikan kita telah
memasangkan bilangan G\"odel dengan $!A$, $x$, dan $t$---misalnya, masing-masing
$k$, $l$, dan $m$. Skema yang sama memasangkan bilangan G\"odel dengan
$\Subst{!A}{t}{x}$, misalnya~$n$. Pemetaan ini---dari $k$, $l$, dan $m$ ke
$n$---merupakan analog aritmetis dari operasi substitusi. Ketika operasi
substitusi memetakan $!A$, $x$, dan $t$ ke $\Subst{!A}{t}{x}$, fungsi
substitusi yang telah diaritmetisasi memetakan bilangan G\"odel $k$, $l$, dan
$m$ ke bilangan G\"odel~$n$. Kita akan melihat bahwa fungsi ini rekursif
primitif.

Aritmetisasi sintaksis bukan sekadar hal yang menarik secara abstrak,
meskipun pada mulanya merupakan wawasan yang tidak sepele bahwa bahasa seperti
bahasa aritmetika, yang tidak dilengkapi mekanisme untuk ``membicarakan''
bahasa, ternyata dapat memformalkan sifat-sifat ekspresi yang kompleks. Dari
sini hanya diperlukan satu langkah kecil untuk menanyakan apa yang dapat
\emph{dibuktikan} oleh suatu teori dalam bahasa ini, seperti aritmetika Peano,
tentang bahasanya sendiri (termasuk, misalnya, apakah !!{sentence}s dapat
dibuktikan atau benar). Hal ini membawa kita pada teorema pembatas G\"odel
yang terkenal (tentang ketakterbuktian) dan teorema Tarski (tentang
ketakterdefinisian kebenaran). Namun, kiat mengaritmetisasi sintaksis juga
penting untuk membuktikan beberapa hasil penting dalam teori komputabilitas,
misalnya mengenai daya komputasional teori atau hubungan antara model-model
komputabilitas yang berbeda. Aritmetisasi sintaksis menjadi pola bagi
aritmetisasi objek dan sifat lain. Sebagai contoh, konfigurasi dan komputasi
(misalnya, komputasi mesin Turing) dapat diaritmetisasi dengan cara serupa.
Hal ini memungkinkan simulasi komputasi dalam satu model (misalnya, mesin
Turing) di dalam model lain (misalnya, fungsi rekursif).

\end{document}
